\documentclass[11pt, letterpaper]{article}
\usepackage{amsmath, amssymb}
\usepackage{hyperref}
\usepackage{booktabs}
\usepackage{enumitem}
\usepackage{xcolor}
\usepackage{geometry}
\usepackage{fancyhdr}
\usepackage{titlesec}
\usepackage{tcolorbox}
\usepackage{parskip}
\geometry{margin=1in}

\definecolor{courseblue}{HTML}{002452}
\definecolor{coursegold}{HTML}{FABD0F}
\definecolor{coursered}{HTML}{B90E31}
\definecolor{lightblue}{HTML}{E8EEF7}

\hypersetup{colorlinks=true,linkcolor=courseblue,urlcolor=coursered}

\tcbuselibrary{skins}
\newtcolorbox{guidancebox}{enhanced,colback=lightblue,colframe=courseblue,
  sharp corners,leftrule=4pt,left=6pt,right=6pt,top=4pt,bottom=4pt}

\pagestyle{fancy}\fancyhf{}
\fancyhead[L]{\small\textcolor{courseblue}{\textbf{CISC 473 --- Deep Learning | Fall 2026}}}
\fancyhead[R]{\small\textcolor{courseblue}{Project Proposal}}
\fancyfoot[C]{\small\thepage}
\renewcommand{\headrulewidth}{0.4pt}

\titleformat{\section}{\large\bfseries\color{courseblue}}{}{0em}{}[\titlerule]
\titleformat{\subsection}{\normalsize\bfseries\color{courseblue}}{}{0em}{}

\begin{document}

{\centering
  {\Large\textbf{CISC 473: Deep Learning --- Project Proposal}}\par\vspace{0.5em}
  {\large\textit{P22: Cross-Architecture Knowledge Distillation}}\par\vspace{0.4em}
  \textbf{Group Members:}\par
  Lee Xuan \quad 20630324 \quad (Role: Project Lead)\par
  Fong Yee Xin \quad 20629927 \quad (Role: Engineering Lead)\par
  \textbf{GitHub Repository:} \url{https://github.com/l33zard/cakd.git}\par\vspace{0.3em}
  \textbf{Date:} \today\par\vspace{0.4em}
  \rule{\linewidth}{0.4pt}
}

\section{Problem Statement and Motivation}

Knowledge distillation (KD) is the process of compressing a large, high-capacity
model---the \emph{teacher}---into a smaller, cheaper \emph{student} model that
preserves much of the teacher's predictive accuracy. Rather than training the
student on hard one-hot labels alone, KD trains it to match the teacher's
softened output distribution, which encodes the teacher's learned similarity
structure over classes~[1]. Subsequent work showed that additional supervision
can be extracted from intermediate representations rather than logits
alone~[2,~5].

The efficiency goal is concrete. Serving large models means routing billions of
queries through energy-hungry data centres, and there is a strong push to move
inference onto personal and embedded hardware---smartphones, laptops, wearables
and automotive chips---so that features work instantly, offline, and without
transmitting sensitive data. Such devices impose strict memory, latency and
power budgets, and typically execute on CPUs or small accelerators where
parameter count and FLOPs directly determine feasibility. Distillation allows a
compact student to inherit teacher behaviour while cutting inference cost,
making real-time deployment in robotics, autonomous systems and edge vision
practical.

\emph{Cross-architecture} KD is the harder setting in which teacher and student
belong to different architectural families. We focus on the gap between
convolutional neural networks (CNNs) and vision transformers (ViTs). CNNs encode
images as hierarchical spatial feature maps produced by local convolutions and
carry a strong locality and translation-equivariance prior; ViTs~[11] encode images
as sequences of patch tokens related by global self-attention. Their intermediate
representations therefore differ in structure, dimensionality and inductive
bias, so feature-level transfer methods designed for CNN-to-CNN pairs do not
apply directly, and logit-only KD discards most of the teacher's internal
knowledge.

This project benchmarks established KD methods on CIFAR-100 and proposes an
intermediate-feature alignment strategy for CNN-to-ViT transfer, evaluated under
a matched parameter and FLOP budget and diagnosed using representational
similarity analysis.

\section{Literature Survey}

\subsection{Related Works}

\paragraph{[1] Hinton, Vinyals \& Dean, 2015. ``Distilling the Knowledge in a Neural Network.'' NeurIPS Deep Learning Workshop. arXiv:1503.02531.}
The paper introduces vanilla KD: the student minimises a weighted combination of
the cross-entropy loss on ground-truth labels and a Kullback--Leibler divergence
against the teacher's logits softened by a temperature $\tau$. The softened
distribution exposes ``dark knowledge''---the relative probabilities the teacher
assigns to incorrect classes---and the authors demonstrate consistent gains for
compact students on MNIST and large-scale speech recognition, as well as for
distilling ensembles into single models. Vanilla KD is our first baseline and
the logit-level loss term that our proposed alignment loss is added on top of;
it is also the only listed method that is trivially architecture-agnostic, since
it operates purely on output distributions.

\paragraph{[2] Tian, Krishnan \& Isola, 2020. ``Contrastive Representation Distillation.'' ICLR. arXiv:1910.10699.}
CRD reformulates distillation as maximising a lower bound on the mutual
information between teacher and student representations, using a contrastive
objective that pulls together teacher--student features of the same input and
pushes apart features of different inputs drawn from a memory bank. It
outperforms vanilla KD and a broad suite of feature-matching methods on
CIFAR-100 and ImageNet, including several teacher--student pairs with differing
CNN families. CRD is a strong feature-level baseline for us; unlike our work, it
compares pooled global embeddings and does not address the token-versus-spatial
mismatch that arises when the student is a transformer.

\paragraph{[3] Zhao, Cui, Song, Qiu \& Liang, 2022. ``Decoupled Knowledge Distillation.'' CVPR. arXiv:2203.08679.}
DKD decomposes the classical KD loss into target-class knowledge distillation
(TCKD) and non-target-class knowledge distillation (NCKD), showing analytically
that the standard formulation couples the NCKD term to the teacher's confidence
on the target class and thereby suppresses exactly the signal that carries most
of the dark knowledge. Decoupling and independently reweighting the two terms
recovers or exceeds state-of-the-art feature-based methods on CIFAR-100,
ImageNet and MS-COCO at essentially no extra training cost. DKD serves as our
strongest logit-only baseline and establishes the reference our cross-architecture
method must beat to justify its additional alignment module.

\paragraph{[4] Chen, Liu, Zhao \& Jia, 2021. ``Distilling Knowledge via Knowledge Review.'' CVPR. arXiv:2104.09044.}
ReviewKD argues that conventional feature distillation pairs layers at the same
depth, whereas a student's deeper layers can profitably ``review'' the teacher's
shallower features. It introduces cross-stage connections together with an
attention-based fusion module and a hierarchical context loss to aggregate
multi-level teacher features, yielding consistent improvements on
classification, detection and segmentation. ReviewKD directly motivates our
study of \emph{which} layer pairs to align: because CNN and ViT depth semantics
do not correspond, layer-pair selection becomes an explicit ablation axis in our
work rather than a fixed design choice.

\paragraph{[5] Romero, Ballas, Kahou, Chassang, Gatta \& Bengio, 2015. ``FitNets: Hints for Thin Deep Nets.'' ICLR. arXiv:1412.6550.}
FitNets is the original intermediate-feature distillation method: a ``hint''
layer in the teacher supervises a ``guided'' layer in the student through an
$L_2$ loss, with a learnable regressor inserted to reconcile the mismatched
feature dimensionalities before the standard KD objective is applied. This made
it possible to train students that are deeper and thinner than their teachers.
Our projection module is a direct descendant of the FitNets regressor,
generalised from a channel-dimension mismatch between two CNNs to a
structural mismatch between a spatial feature map and a patch-token sequence.

\paragraph{[6] Kornblith, Norouzi, Lee \& Hinton, 2019. ``Similarity of Neural Network Representations Revisited.'' ICML. arXiv:1905.00414.}
The paper analyses measures of representational similarity and shows that
Centered Kernel Alignment (CKA) is invariant to orthogonal transformation and
isotropic scaling while remaining sensitive to genuine differences, unlike CCA
variants; this makes it reliable for comparing layers of different widths or of
independently trained networks. We adopt CKA both as a diagnostic---quantifying
how CNN-teacher and ViT-student representations relate before and after
distillation---and, optionally, as a differentiable alignment objective.

\paragraph{[7] Park, Kim, Lu \& Cho, 2019. ``Relational Knowledge Distillation.'' CVPR. arXiv:1904.05068.}
RKD transfers the \emph{structure} of the teacher's embedding space rather than
individual activations, penalising discrepancies in pairwise distances and
triplet angles among samples in a mini-batch. Because it constrains only
relations between examples, it is insensitive to the absolute geometry and
dimensionality of either representation space. This property is attractive for
cross-architecture transfer and gives us a natural secondary comparison for our
projection-based alignment.

\subsection{Open Problem}

Existing feature-level distillation methods [2,~4,~5,~7] are developed and
evaluated almost exclusively on CNN-to-CNN pairs, while the methods that
transfer cleanly across architectures [1,~3] use logits only and therefore
discard intermediate knowledge. Cross-architecture KD with a CNN teacher and a
ViT student consequently lacks (i) a systematic study of how to align spatial
feature maps with patch-token sequences, including which layer pairs should be
matched when depth semantics do not correspond; (ii) an efficiency analysis that
reports accuracy jointly with parameters, FLOPs and training cost under a
matched budget, as required for CPU and edge deployment; and (iii) a principled
diagnosis of \emph{why} transfer succeeds or fails---temperature and
loss-weighting ablations are rarely reported for this setting, and CKA-based
analysis of teacher--student representational similarity [6] has not been applied
to it. This project addresses all three gaps.

\section{Proposed Methodology}

\subsection{Baseline}

\textbf{Dataset.} All experiments use CIFAR-100~[12]: 60{,}000 $32\times32$ colour
images across 100 classes, with the standard 50{,}000/10{,}000 train/test split.
We hold out 5{,}000 training images as a validation set for hyperparameter and
checkpoint selection. Standard augmentation (random crop with 4-pixel padding,
horizontal flip, per-channel normalisation) is applied. If compute permits, we
will scale the best configuration to TinyImageNet ($64\times64$, 200 classes) to
test whether conclusions hold at higher resolution.

\textbf{Models.} The CNN-to-CNN teacher is ResNet-110~[8] trained on CIFAR-100 with
the standard recipe; the corresponding student is MobileNetV2~[9], giving a
conventional compression baseline. For the cross-architecture setting we use a
ResNet-50~[8] teacher, whose wider stage outputs are better matched to a token
embedding dimension, and a DeiT-Tiny~[10] student (192-dimensional embedding, 12
layers, 3 heads) adapted to $32\times32$ inputs with a reduced patch size.

\textbf{Reproduced baselines.} We will reproduce Vanilla KD~[1], CRD~[2],
DKD~[3] and ReviewKD~[4] starting from the RepDistiller framework and the
authors' official releases, retaining their published hyperparameters where
available so that reported numbers can be validated before we modify anything.

\textbf{Metrics.} Top-1 test accuracy is the primary metric; we additionally
report top-5 accuracy, parameter count, inference FLOPs and wall-clock training
time per run, so that any accuracy gain can be read against its cost.

\subsection{Proposed Extension}

CNN feature maps $F_T \in \mathbb{R}^{C \times H \times W}$ and ViT token
sequences $F_S \in \mathbb{R}^{N \times D}$ cannot be compared directly. We
propose a lightweight intermediate-feature alignment module that operates
entirely outside the student, leaving the DeiT architecture unmodified. The
teacher map is spatially pooled or bilinearly interpolated to a $\sqrt{N} \times
\sqrt{N}$ grid, flattened into $N$ vectors, and passed through a learnable
$1\times1$ projection $\mathbb{R}^{C} \to \mathbb{R}^{D}$. The resulting
pseudo-tokens are aligned with the student's tokens by an alignment loss
$\mathcal{L}_{\text{align}}$, instantiated either as a differentiable CKA
objective~[6] or as an attention-map transfer loss that matches the teacher's
spatial activation statistics to the student's self-attention maps. The total
objective is
\[
\mathcal{L} = (1-\alpha)\,\mathcal{L}_{\text{CE}}
            + \alpha\,\tau^{2}\,\mathcal{L}_{\text{KD}}(\tau)
            + \beta\,\mathcal{L}_{\text{align}} .
\]
Because the projection is discarded at inference, the deployed student's
parameter count and FLOPs are unchanged.

\textbf{Hypothesis.} Explicit CKA or attention-based alignment of intermediate
representations gives the DeiT-Tiny student supervision that logit-only KD
cannot provide, improving top-1 accuracy over Vanilla KD, CRD, DKD and ReviewKD
at a matched parameter and FLOP budget.

\textbf{Ablations.} We will sweep the temperature $\tau \in \{1,2,4,8,16\}$, the
distillation weight $\alpha$, the alignment weight $\beta$, the choice of
alignment loss, and the teacher--student layer pairs (early, middle, late and
multi-layer combinations), following the layer-pairing question raised by
ReviewKD~[4].

\subsection{Evaluation Plan}

\begin{itemize}[noitemsep]
  \item \textbf{Dataset and splits.} CIFAR-100 with the standard train/test
        split and a fixed 5{,}000-image validation subset held out from training;
        identical splits, augmentation and schedules across all methods.
        Optional TinyImageNet scale-up.
  \item \textbf{Baselines.} DeiT-Tiny and MobileNetV2 trained from scratch
        (lower bound), Vanilla KD~[1], CRD~[2], DKD~[3], ReviewKD~[4], and the
        teacher's own accuracy (upper bound).
  \item \textbf{Metrics.} Top-1 and top-5 accuracy, parameters, FLOPs, training
        time, and layer-wise CKA similarity between teacher and student measured
        before and after distillation.
  \item \textbf{Statistical reliability.} Every main result uses at least three
        independent runs with fixed seeds, reported as mean $\pm$ standard
        deviation; ablation sweeps use a single seed for screening, with the
        selected configurations rerun over three seeds.
\end{itemize}

\subsection{Early Novel Contribution (Optional Bonus)}

We propose two concrete contributions. First, a systematic CKA-based study of
CNN-to-ViT knowledge distillation: layer-wise CKA heatmaps between a ResNet
teacher and a DeiT-Tiny student, measured before and after distillation under
each baseline method, which to our knowledge has not been reported and would
give a diagnostic account of where cross-architecture transfer breaks down.
Second, a lightweight intermediate-feature alignment module that requires no
modification to the ViT architecture and adds no inference-time cost, since the
projection is used only during training. If the ablations show that alignment
improves student accuracy under a matched parameter and FLOP budget, and that
the improvement is accompanied by measurably higher CKA similarity, we will
submit this as our +2 bonus contribution.

\section{Project Timeline and Roles}

\begin{center}
\begin{tabular}{lp{7cm}l}
\toprule
\textbf{Week} & \textbf{Tasks} & \textbf{Responsible} \\
\midrule
1--2 & Read DKD, CRD and ReviewKD; environment and repository setup; CIFAR-100 data pipeline; train ResNet-110 teacher and from-scratch student baselines; \textbf{Proposal due} & Lee Xuan \& Fong Yee Xin \\
3 & Vanilla KD implementation; temperature $\tau$ and distillation-weight $\alpha$ ablations & Lee Xuan \\
4 & CRD distillation baseline; memory-bank and negative-sampling configuration & Fong Yee Xin \\
5 & DKD distillation baseline; ResNet-50 and DeiT-Tiny setup for cross-architecture ResNet-50 $\rightarrow$ DeiT-Tiny transfer & Fong Yee Xin \\
6 & ReviewKD baseline; layer-wise CKA analysis before and after distillation; \textbf{Midterm due} & Lee Xuan \\
7 & Intermediate-feature alignment module for the CNN $\rightarrow$ ViT case; projection and alignment-loss implementation & Fong Yee Xin \\
8 & Extended experiments and ablations ($\beta$, layer pairs, alignment loss); efficiency analysis of accuracy vs.\ FLOPs & Lee Xuan \& Fong Yee Xin \\
9 & Results analysis; accuracy tables with mean $\pm$ std; CKA heatmaps; report writing & Lee Xuan \\
10 & Final report, reproducibility checks, README and code release; \textbf{Final due} & Lee Xuan \& Fong Yee Xin \\
\bottomrule
\end{tabular}
\end{center}

\subsection{Contribution Declaration}

\begin{center}
\begin{tabular}{lp{7cm}c}
\toprule
\textbf{Member} & \textbf{Primary Responsibilities} & \textbf{Estimated \%} \\
\midrule
Lee Xuan (Project Lead) & Project coordination, literature review, Vanilla KD and ReviewKD baselines, CKA analysis, efficiency analysis, report writing & 50\% \\
Fong Yee Xin (Engineering Lead) & Data and training pipeline, CRD and DKD baselines, cross-architecture setup, feature-alignment module implementation, ablation design and execution, evaluation code and README & 50\% \\
\bottomrule
\end{tabular}
\end{center}

\section{References}

\begin{enumerate}[label={[\arabic*]},leftmargin=*]
  \item Hinton, G., Vinyals, O., and Dean, J., 2015. ``Distilling the Knowledge in a Neural Network.'' \textit{NeurIPS Deep Learning and Representation Learning Workshop}. arXiv:1503.02531.
  \item Tian, Y., Krishnan, D., and Isola, P., 2020. ``Contrastive Representation Distillation.'' \textit{International Conference on Learning Representations (ICLR)}. arXiv:1910.10699.
  \item Zhao, B., Cui, Q., Song, R., Qiu, Y., and Liang, J., 2022. ``Decoupled Knowledge Distillation.'' \textit{IEEE/CVF Conference on Computer Vision and Pattern Recognition (CVPR)}. arXiv:2203.08679.
  \item Chen, P., Liu, S., Zhao, H., and Jia, J., 2021. ``Distilling Knowledge via Knowledge Review.'' \textit{IEEE/CVF Conference on Computer Vision and Pattern Recognition (CVPR)}. arXiv:2104.09044.
  \item Romero, A., Ballas, N., Kahou, S. E., Chassang, A., Gatta, C., and Bengio, Y., 2015. ``FitNets: Hints for Thin Deep Nets.'' \textit{International Conference on Learning Representations (ICLR)}. arXiv:1412.6550.
  \item Kornblith, S., Norouzi, M., Lee, H., and Hinton, G., 2019. ``Similarity of Neural Network Representations Revisited.'' \textit{International Conference on Machine Learning (ICML)}. arXiv:1905.00414.
  \item Park, W., Kim, D., Lu, Y., and Cho, M., 2019. ``Relational Knowledge Distillation.'' \textit{IEEE/CVF Conference on Computer Vision and Pattern Recognition (CVPR)}. arXiv:1904.05068.
  \item He, K., Zhang, X., Ren, S., and Sun, J., 2016. ``Deep Residual Learning for Image Recognition.'' \textit{IEEE Conference on Computer Vision and Pattern Recognition (CVPR)}. arXiv:1512.03385.
  \item Sandler, M., Howard, A., Zhu, M., Zhmoginov, A., and Chen, L.-C., 2018. ``MobileNetV2: Inverted Residuals and Linear Bottlenecks.'' \textit{IEEE/CVF Conference on Computer Vision and Pattern Recognition (CVPR)}. arXiv:1801.04381.
  \item Touvron, H., Cord, M., Douze, M., Massa, F., Sablayrolles, A., and J\'egou, H., 2021. ``Training Data-Efficient Image Transformers \& Distillation Through Attention.'' \textit{International Conference on Machine Learning (ICML)}. arXiv:2012.12877.
  \item Dosovitskiy, A., Beyer, L., Kolesnikov, A., Weissenborn, D., Zhai, X., Unterthiner, T., Dehghani, M., Minderer, M., Heigold, G., Gelly, S., Uszkoreit, J., and Houlsby, N., 2021. ``An Image Is Worth 16x16 Words: Transformers for Image Recognition at Scale.'' \textit{International Conference on Learning Representations (ICLR)}. arXiv:2010.11929.
  \item Krizhevsky, A., 2009. ``Learning Multiple Layers of Features from Tiny Images.'' Technical Report, University of Toronto.
\end{enumerate}

\end{document}
